\documentclass[12pt]{amsart}
\usepackage{amsthm}
\usepackage{url}
\newtheorem{lemma}{\bf Lemma}[section]
\newtheorem{proposition}[lemma]{\bf Proposition}
\newtheorem{theorem}[lemma]{\bf Theorem}
\newtheorem{observation}[lemma]{\bf Observation}
\newtheorem{definition}[lemma]{\bf Definition}
%------------
\gdef\cal{\mathcal}
\gdef\frak{\mathfrak}
\gdef\a{\omega}
%------------
\parskip = 2mm
\parindent = 0mm

\begin{document}
\title{${\Omega \choose \Gamma}\text{ and the }\delta\text{-property}$}
\author{Dominic van der Zypen}
\address{Federal Office of Borders and Customs Security,
Taubenstrasse 16, CH-3003 Bern, Switzerland}
\email{dominic.zypen@gmail.com}

\begin{abstract}
Gerlits and Nagy \cite{gamma} introduced 
the selection principle ${\Omega\choose \Gamma}$ as well
as a weaker limit concept called the $\delta$-property, and asked
if they are equivalent. We give a thorough introduction
to the concepts needed to understand this question and provide
a positive answer.
\end{abstract}

\maketitle
%-----------------
\section{Set limits and covers}
% -  -  -  -  -  -
\subsection{Set limits} For a sequence of subsets $(X_n)_{n\in\mathbb{N}}$
of a set $X$, we define its {\em (lower) limit} by setting 
$$\liminf X_n = \bigcup_m\bigcap_{n\geq m}X_n.$$

% -  -  -  -  -  -
\subsection{Covers of a topological space}
Throughout this note, let $(X,\tau)$ be a non-empty topological space.
\begin{definition} \label{covdef}
 We say that ${\cal U} \subseteq \tau$ is an {\em open cover}, or
  {\em cover} for short, if \begin{enumerate}
    \item $X\notin {\cal U}$,
    \item $\bigcup {\cal U} = X$.
  \end{enumerate}
\end{definition}
For $x\in X$ we set ${\cal U}_x = \{U \in {\cal U}:
x\in U\}$ and call ${\cal U}_x$ the
{\em star} of $x$ (with respect to ${\cal U}$).
\begin{definition}
	We call a cover ${\cal U}\subseteq (\tau\setminus\{X\})$
\begin{enumerate}
  %\item {\em large} if ${\cal U}_x$ is infinite for every $x\in X$,
  \item an $\omega$-cover if every finite subset of $X$ is contained
    in some member of ${\cal U}$,
  \item a $\gamma$-cover if ${\cal U}$ is infinite, and
    for every $x\in X$ the set ${\cal U}
    \setminus {\cal U}_x$ is finite.
\end{enumerate}
\end{definition}
Let $\Omega$ and $\Gamma$ denote the collections
of (open) $\omega$-covers and $\gamma$-covers, respectively. 

%  -  -  -  -  -  -
\subsection{The selection principle} Let ${\mathfrak U}, {\mathfrak V}$ be
families of covers of $X$. We define the property ${{\frak U}\choose
{\frak V}}$, read ``$\frak U$ choose $\frak V$'', as follows:

\begin{quote}
  ${{\frak U}\choose {\frak V}}\; :$ For each ${\cal U}\in {\frak U}$
  there is ${\cal V}\subseteq {\cal U}$ such that ${\cal V}\in {\frak V}$.
\end{quote}

\subsection{The $\delta$-property} We say that a topological
space $(X,\tau)$ has the $\delta$-{\em property} if 
for every $\omega$-cover ${\cal U}$ of $X$ there is a sequence
$(U_n)_{n\in\mathbb{N}}$ of members of ${\cal U}$ such that 
$$X = \liminf U_n = \bigcup_m\bigcap_{n\geq m} U_n.$$

%-----------------
\section{Main result}
Whenever a space $X$ satisfies ${\Omega\choose \Gamma}$, it has
the $\delta$-property (see \cite{gamma} and \cite{spm}), 
so the $\delta$-property
implies ${\Omega\choose\Gamma}$. This gives rise
to the 

{\bf Question} (Gerlits-Nagi \cite{gamma}). Is the $\delta$-property
equivalent to ${\Omega\choose\Gamma}$?

%- - - - - - -
\subsection{Strategy} If we have a space $X$ with the $\delta$-property,
and an $\omega$-cover ${\cal U}$, we get a sequence 
$(U_n)_{n\in\mathbb{N}}$ of members of ${\cal U}$ 
such that $$\liminf U_n = \bigcup_m\bigcap_{n\geq m}U_n = X.$$ 
We want to find
a subcover ${\cal U}^*\subseteq {\cal U}$ such that 
${\cal U}^*$ is a $\gamma$-cover of $X$, with the aim
to establish ${\Omega\choose \Gamma}$. The idea is to pick 
$${\cal U}^* = \{U_n:n\in\mathbb{N}\}.$$
In order to prove that ${\cal U}^*$ is a $\gamma$-cover, we need to show
that 
\begin{enumerate}
\item ${\cal U}^*$ is a cover, 
\item ${\cal U}^*$ is infinite, and
\item ${\cal U}^*\setminus {\cal U}^*_x$ is finite for all $x\in X$.
\end{enumerate}
Then we can conclude that every $\omega$-cover contains a $\gamma$-cover
and conclude that ${\Omega\choose\Gamma}$ holds.

The next two lemmas take care of tasks (1) and (2).

\begin{lemma}\label{omegalemma}
	If $(X,\tau)$ is a topological space having an $\omega$-cover 
	${\cal W}$ (that is, ${\cal W}\in \Omega$), 
	then the base set $X$ is infinite.
\end{lemma}
{\em Proof.} If $X$ were finite, then since ${\cal W}$ is an $\omega$-cover,
there is $W_0\in {\cal W}$ with $X\subseteq W_0$. As $V\subseteq X$ for every
$V\in {\cal W}$, we get $W_0 = X\in {\cal W}$, contradicting the definition of a cover
as in \ref{covdef}. \hfill{$\Box$}

\begin{lemma}\label{gammalemma}
Let $(X,\tau)$ be a space with an $\omega$-cover ${\cal W}$.
Suppose that there is a sequence $(U_n)_{n\in\mathbb{N}}$
of elements of ${\cal W}$ such that $X =\liminf U_n$.
Set ${\cal U}^* := \{U_n:n\in \omega\}$. Then
\begin{enumerate}
\item ${\cal U}^*$ is a cover, 
\item ${\cal U}^*$ is infinite.
\end{enumerate}
\end{lemma}
{\em Proof.} (1) This follows directly from $X = \bigcup_m\bigcap_{n\geq m}U_n$.

(2) 
Recall that we know from lemma \ref{omegalemma} that the base set
$X$ itself is infinite. 

Suppose that ${\cal U}^* := \{U_n:n\in \omega\}$ were finite.
That means that ${\cal U}^*=\{U_{n_0},U_{n_1},\ldots, U_{n_k}\}$ for some
indices $n_0,\ldots,n_k$. There is at least one $j\in\{0,\ldots,k\}$ such that 
$U_{n_j}$ appears infinitely often in the sequence $(U_n)_{n\in\mathbb{N}}$. 
By the definition of a cover (\ref{covdef}), we have $U_{n_j} \neq X$, so pick 
$x_0\in X\setminus U_{n_j}$. As
$U_{n_j}$ appears infinitely often in the sequence $(U_n)_{n\in\mathbb{N}}$,
we get $$x_0\notin \bigcup_m\bigcap_{n\geq m}U_n=\liminf U_n,$$
contradicting $X = \liminf U_n$. So ${\cal U}^*$ is infinite.\hfill{$\Box$}

%\/\/\/\/\/\/\/\/\/\/\/\/\/\/\/\/\/\/
\begin{theorem}\label{main}
Let $(X,\tau)$ be a topological space. Then $X$ has
the $\delta$-property if and only if it 
satisfies ${\Omega\choose\Gamma}$.
\end{theorem}
%\/\/\/\/\/\/\/\/\/\/\/\/\/\/\/\/\/\/
{\it Proof.}

$\Leftarrow:$ Established in \cite{gamma}.

$\Rightarrow:$ 

{\sl Case 1.} $\Omega = \emptyset$. If $(X,\tau)$ has no $\omega$-cover,
then ${\Omega\choose\Gamma}$ holds vacuously.

{\sl Case 2:} $\Omega\neq\emptyset$. 
Let ${\cal U}$ be an open $\omega$-cover of $(X,\tau)$.
By assumption $(X,\tau)$ has the $\delta$-property, so 
there is a sequence $(U_n)_{n\in\mathbb{N}}$ of members of ${\cal U}$
such that $$X = \liminf U_n=\bigcup_m\bigcap_{n\geq m}U_n.$$ 
This means that
\begin{quote}
$(\star)$ for every $x\in X$ there is $m\in \mathbb{N}$ such that
$x\in U_n$ for all $n\geq m$.
\end{quote}
Let ${\cal U}^* = \{U_n: n\in\mathbb{N}\}$. 

We want to show that ${\cal U}^*$ is a $\gamma$-cover.

First, by lemma \ref{gammalemma}, we know that ${\cal U}^*$ 
is a cover of $X$, and moreover the same lemma establishes
that ${\cal U}^*$ is infinite.

Second, we want to establish that ${\cal U}^*\setminus {\cal U}^*_x$
is finite for all $x\in X$.  Let $x\in X$ be an arbitrary
element. By $(\star)$, there is $m\in \mathbb{N}$
such that $x\in U_n$ for all $n\geq m$. Therefore $${\cal U}^*_x \supseteq
\{U_n: n\geq m\}.$$ This implies
$$\big({\cal U}^*\setminus {\cal U}^*_x\big) \subseteq \big\{U_0, U_1,\ldots, 
U_{m}\big\},$$ therefore ${\cal U}^*\setminus {\cal U}^*_x$ is finite.

So we have shown that ${\cal U}^*$ is a $\gamma$-cover.

Taking stock, for any $\omega$-cover ${\cal U}$ we found a $\gamma$-cover
${\cal U}^*\subseteq {\cal U}$, and we conclude that the property 
${\Omega\choose\Gamma}$ holds.
\hfill{$\Box$}

{\bf AI disclosure.} AI has not been used at any stage of the writing
of this note.

%-------------------
\begin{thebibliography}{9}
  \bibitem{gamma} J.~Gerlits and Zs.~Nagy, {\it Some properties of
    $C(X)$, I}, Topology and its Applications {\bf 14} (1982), 151–161.
  \bibitem{sakai} M.~Sakai, {\it Special subsets of reals characterizing
	  local properties of function spaces} in {\bf Selection Principles
		and Covering Properties in Topology} (L.~Ko\v{c}inac, ed.)
		Quaderni di Matematica {\bf 18} (2007), 195-225.
  \bibitem{spm} Boaz Tsaban, {\it SPM (Selection Principles in Mathematics)
	  bulletin No.~23} (2008), \url{https://arxiv.org/pdf/0801.3181}
\end{thebibliography}
\end{document}
